\begin{trapbox}
\begin{center}
\begin{tikzpicture}[font=\small,>=Stealth,line cap=round,line join=round]
\path[use as bounding box] (-6.5,-3.55) rectangle (6.5,3.55);
\draw[black!20,rounded corners=3pt] (-6.5,-3.55) rectangle (6.5,3.55);
\node[font=\small\bfseries] at (0,3.08) {同样的渐近线，不能保证同样的对称轴};

\foreach \dx/\mode/\ttl in {-3.25/0/{$y=1/x$},3.25/1/{$y=1/x+1/x^2$}} {
\begin{scope}[shift={(\dx,0)},x=.65cm,y=.65cm]
\draw[->,black!55,dashed] (-3.6,0)--(3.6,0) node[right] {$x$};
\draw[->,black!55,dashed] (0,-3.6)--(0,3.6) node[above] {$y$};
\draw[CProof,dotted] (-3,-3)--(1.8,1.8);
\node[fill=white] at (-2.4,-2.8) {$y=x$};
\begin{scope}
\clip (-3.3,-3.3) rectangle (3.3,3.3);
\draw[CExplanation,very thick,domain=.25:3.3,samples=100] plot (\x,{1/\x+\mode/(\x*\x)});
\draw[CExplanation,very thick,domain=-3.3:-.25,samples=100] plot (\x,{1/\x+\mode/(\x*\x)});
\end{scope}
\node[fill=white] at (0,3.95) {\ttl};
\ifnum\mode=0
\fill[CExplanation] (2,.5) circle(2pt) node[above right] {$(2,\frac12)$};
\fill[CExplanation] (.5,2) circle(2pt) node[above left] {$(\frac12,2)$};
\draw[CProof,<->] (2,.5)--(.5,2);
\else
\fill[CExplanation] (1,2) circle(2pt) node[above right] {$P(1,2)$};
\draw[CTrap,->] (1,2)--(2,1);
\draw[CTrap,thick,fill=white] (2,1) circle(2.6pt);
\node[right] at (2.1,1.5) {$P'(2,1)$};
\fill[CExplanation] (2,.75) circle(1.6pt);
\node at (2,-1.5) {$f(2)=\frac34\ne1$};
\fi
\end{scope}}

\node[align=center] at (0,-3.05) {两图同尺度；右图 $P'$ 落在曲线外，一点即可排除 $y=x$。};
\end{tikzpicture}
\end{center}


\begin{description}[style=nextline, leftmargin=2.8em, labelsep=0.5em, itemsep=0.35em]
  \item[\textbf{\textcolor{CTrap}{易错点一（把几个对称点当成整条曲线对称）：}}] 找到曲线上两三个点照镜子以后还在曲线上，就直接宣称有对称轴。
  \item[\textbf{\textcolor{CTrap}{正确理解：}}] 证明轴对称必须对\textbf{任意点}成立；否定对称只需\textbf{一个反例点}。例如 $y=1/x+1/x^2$ 的 $(1,2)$ 足以否定 $y=\pm x$。
  \item[\textbf{\textcolor{CTrap}{错因分析：}}] 把“存在某个点符合”错误当成“所有点都符合”。
\end{description}

\medskip
\begin{description}[style=nextline, leftmargin=2.8em, labelsep=0.5em, itemsep=0.35em]
  \item[\textbf{\textcolor{CTrap}{易错点二（有两条渐近线就认定角平分线是轴）：}}] 看到渐近线 $x=0,y=0$ 就立即断言曲线关于 $y=\pm x$ 对称。
  \item[\textbf{\textcolor{CTrap}{正确理解：}}] 对\textbf{双曲线}可以这样寻找轴；一般曲线不能这样推出。反例 $y=1/x+1/x^2$ 有相同两条渐近线，却不关于 $y=\pm x$ 对称。
  \item[\textbf{\textcolor{CTrap}{错因分析：}}] 渐近线描述曲线向远处延伸时的方向，不能决定曲线上所有点的位置。
\end{description}

\medskip
\begin{description}[style=nextline, leftmargin=2.8em, labelsep=0.5em, itemsep=0.35em]
  \item[\textbf{\textcolor{CTrap}{易错点三（混淆 $y=x$ 与 $y=-x$）：}}] 关于两条斜线对称都只交换 $x,y$，或者都只换符号。
  \item[\textbf{\textcolor{CTrap}{正确理解：}}] 关于 $y=x$：$(x,y)\to(y,x)$；关于 $y=-x$：$(x,y)\to(-y,-x)$。关于原点中心对称才是 $(x,y)\to(-x,-y)$。
  \item[\textbf{\textcolor{CTrap}{错因分析：}}] 对 $P(2,\frac12)$，前文两幅坐标交换图的镜像分别是 $(\frac12,2)$、$(-\frac12,-2)$；关于原点的中心对称点则是 $(-2,-\frac12)$。三者位置不同，不能混用。
\end{description}

\medskip
\begin{center}
\begin{tikzpicture}[font=\small,>=Stealth,line cap=round,line join=round]
\path[use as bounding box] (-6.5,-3.55) rectangle (6.5,3.55);
\draw[black!20,rounded corners=3pt] (-6.5,-3.55) rectangle (6.5,3.55);
\node[font=\small\bfseries] at (0,3.08) {定义域也要对称：只留半支就会缺少镜像};

\foreach \dx/\mode/\ttl in {-3.25/0/{完整抛物线},3.25/1/{只取 $x\ge2$}} {
\begin{scope}[shift={(\dx,-1.8)},x=.9cm,y=.85cm]
\draw[->,black!40] (-2.5,0)--(2.5,0);
\draw[CProof,dashed] (0,0)--(0,4.2);
\node[right=3pt] at (0,3.9) {$x=2$};
\ifnum\mode=0
\draw[CExplanation,very thick,domain=-1.65:1.65,samples=80] plot (\x,{\x*\x+1});
\fill[CExplanation] (-1,2) circle(2pt);
\else
\draw[black!25,dotted,domain=-1.65:0,samples=40] plot (\x,{\x*\x+1});
\draw[CExplanation,very thick,domain=0:1.65,samples=40] plot (\x,{\x*\x+1});
\draw[CTrap,fill=white,thick] (-1,2) circle(3pt);
\fi
\fill[CExplanation] (1,2) circle(2pt);
\draw[<->,CProof] (-1,2)--(1,2);
\node[below] at (-1,0) {$1$}; \node[below] at (1,0) {$3$};
\node at (0,4.9) {\ttl};
\end{scope}}

\node[align=center] at (0,-3.05) {公式同为 $y=(x-2)^2+1$；右图的 $(1,2)$ 不属于实际图像。};
\end{tikzpicture}
\end{center}


\begin{description}[style=nextline, leftmargin=2.8em, labelsep=0.5em, itemsep=0.35em]
  \item[\textbf{\textcolor{CTrap}{易错点四（只代换表达式，不检查定义域）：}}] 只算出 $f(h+t)=f(h-t)$ 的形式相同，却忘记一个自变量可能不在定义域内。
  \item[\textbf{\textcolor{CTrap}{正确理解：}}] 对于函数图像，$x\in D$ 当且仅当 $2h-x\in D$，并且成对的函数值相等。分式需检查分母不为零；只取半支曲线时不能擅自补上另一半。
  \item[\textbf{\textcolor{CTrap}{错因分析：}}] 轴对称判断的是\textbf{实际存在的点的集合}，不是只看展开后的符号。
\end{description}

\medskip
\begin{description}[style=nextline, leftmargin=2.8em, labelsep=0.5em, itemsep=0.35em]
  \item[\textbf{\textcolor{CTrap}{易错点五（认为只有方程完全一样才算对称）：}}] 把 $F(x',y')=F(x,y)$ 当成对称的必要条件。
  \item[\textbf{\textcolor{CTrap}{正确理解：}}] 必要且充分的是 $F(x,y)=0\Longleftrightarrow F(x',y')=0$。例如 $F=(x^2+y^2-1)(x^2+2)$，表示的仍是单位圆；关于 $y=x$ 交换后因第二因子不同，两个\emph{表达式}通常不等，但两方程的解集一样。
  \item[\textbf{\textcolor{CTrap}{错因分析：}}] 忽视“方程的表达式”和“方程表示的曲线”是两个不同概念。
\end{description}

\medskip
\begin{description}[style=nextline, leftmargin=2.8em, labelsep=0.5em, itemsep=0.35em]
  \item[\textbf{\textcolor{CTrap}{易错点六（以为一条曲线一定只有零条、一条或两条轴）：}}] 看到正比例函数，就照抛物线思路只找出一条轴。
  \item[\textbf{\textcolor{CTrap}{正确理解：}}] 一条\textbf{完整直线}关于自身对称，也关于任一条与它垂直且相交的直线对称，所以有无数条对称轴。标准圆也有无数条经过圆心的对称轴。
  \item[\textbf{\textcolor{CTrap}{错因分析：}}] 不同图形的对称轴数量不能机械类比。
\end{description}

\medskip
\begin{description}[style=nextline, leftmargin=2.8em, labelsep=0.5em, itemsep=0.35em]
  \item[\textbf{\textcolor{CTrap}{易错点七（套法向量时漏掉常数或分母）：}}] 对 $Ax+By+C=0$ 直接写 $(x',y')=(x-2A L,y-2B L)$，漏掉 $A^2+B^2$。
  \item[\textbf{\textcolor{CTrap}{正确理解：}}] 先算 $s=(Ax+By+C)/(A^2+B^2)$，再写 $(x',y')=(x-2As,y-2Bs)$。只有 $A^2+B^2=1$ 时才可省去分母。
  \item[\textbf{\textcolor{CTrap}{错因分析：}}] 法线方向 $(A,B)$ 未必长度为 $1$，不能把方程代入值直接当作垂直距离。
\end{description}
\end{trapbox}
